\section{Goldbach's Other Conjecture --- Problem 46}
\label{prob:46}
\noindent\textbf{Source:} \url{https://github.com/KO0I/projectEuler/blob/reorganize/solutions-by-language/solutions/c/goldbachs_other_conjecture_46.c}

\noindent\textbf{Assessment:} Conceptual search, unsafe bounds and missing declarations.

\subsection*{Mathematical goal}
The conjecture under test is that every odd composite integer \(z\) can be written \(z=p+2m^2\) for a prime \(p\) and positive integer \(m\). Find the least counterexample.

\subsection*{How the C code works}
The program reads prime numbers from \verb|million_primes|, generates odd composite candidates using the external helper \verb|isPrime|, then searches pairs \((p,m)\) for the required decomposition. It calls external \verb|powi| to square \verb|m|, advances to the next candidate when a representation is found, and times the search.

\subsection*{Correctness and limitations}
The source does not declare its external helper prototypes or link them directly; it must be built with \verb|support/c/functions/isPrime.c| and \verb|powi.c|. It reads without guarding the million-element array capacity and loops using \verb|i<=1000000|, an out-of-range index. Advancing \verb|i| inside nested loops can also invalidate the current candidate. The expected counterexample is 5777; the current implementation cannot be regarded as a safe computation of it.


\subsection*{Correction to historical C}
{\color{red}\textbf{Correction:} The corrected odd-composite search tests prime-plus-twice-a-square representations directly, eliminating unsafe prime-file indexing and nested-loop index mutations.}
